% Minimal preamble for Open Logic Project sources (pinned commit 1e960beff9ed7835bf3e3f1335e21af3439cd107).
%
% OLP's own master preamble cannot run on the project host: TeX Live 2025 there lacks
% prooftree.sty, ntheorem.sty, extarrows.sty, mathdots.sty, isomath.sty and subfiles.sty.
% Following the project's own rule ("if the author preamble needs packages the toolchain lacks,
% do not invent macros -- use a minimal working preamble"), this file declares, verbatim in
% spirit and numbering behaviour, exactly the theorem-like environments OLP defines in
% sty/open-logic-book-envs.sty (thm/lem/prop/cor/defn/ex/prob/rem/note/case/axiom/conv),
% plus the shorthand accents OLP uses (`!!a`, `!!b`, `!!{...}`).
%
% Author text is untouched; only the packaging is ours.
\documentclass[10pt,a4paper]{article}

\usepackage{amsmath}
\usepackage{amssymb}
\usepackage{amsthm}
\usepackage{graphicx}
\usepackage{hyperref}
\usepackage{enumitem}
\usepackage[framemethod=default]{mdframed}

% ---- OLP shorthand accents (from sty/open-logic.sty) ----
% NB: \! is taken by LaTeX (negative thin space), so we only define the !!a / !!b families
% and the \olup helper used by them.
\newcommand*{\olup}[1]{\MakeUppercase{#1}}
\def\!!!a#1{\olup{#1}}
\def\!!a#1{\olup{#1}}
\def\!!b#1{\olup{#1}}

% ---- theorem-like environments, same family numbering as OLP ----
\theoremstyle{plain}
\newtheorem{thm}{Theorem}[section]
\newtheorem{lem}[thm]{Lemma}
\newtheorem{prop}[thm]{Proposition}
\newtheorem{cor}[thm]{Corollary}
\newtheorem{axiom}{Axiom}

\theoremstyle{definition}
\newtheorem{defn}[thm]{Definition}
\newtheorem{ex}[thm]{Example}
\newtheorem{prob}{Problem}[section]
\newtheorem{probd}{Problem}
\newtheorem{conv}[thm]{Convention}

\theoremstyle{remark}
\newtheorem{rem}[thm]{Remark}
\newtheorem{note}[thm]{Note}
\newtheorem{case}[thm]{Case}

% ---- environments OLP defines as wrappers ----
\newenvironment{probtag}[1]{\begin{prob}}{\end{prob}}
\newenvironment{explain}{}{}
\newenvironment{explain*}{}{}
\newenvironment{editorial}{}{}
\newenvironment{digress}{}{}
\newenvironment{digress*}{}{}
\newenvironment{history}{}{}
\newenvironment{reference}{}{}
\newcommand*{\olfileid}[3]{}
\newcommand*{\olsection}[1]{\section*{#1}}
\newcommand*{\olchapter}[3]{}
\newcommand*{\olpart}[2]{}
\newcommand*{\olimport}[1]{}
\newcommand*{\olinclude}[1]{}
\newcommand*{\olpath}[1]{}
\newcommand*{\ollangid}[1]{}
\newcommand*{\problemtag}[1]{}
\newcommand*{\problemauthor}[1]{}
\newcommand*{\ollabel}[1]{\label{#1}}
% \olref takes 0-3 OPTIONAL arguments plus one mandatory label:
%     \olref{lab}  |  \olref[from]{lab}  |  \olref[from][to]{lab}  |  \olref[a][b][c]{lab}
% A single-argument definition silently swallows "[sth]" and emits a broken reference
% (measured on the host: the text came out as "??iso]x"). The chain below consumes the
% brackets and cross-references only the mandatory argument.
\makeatletter
\newcommand{\olref}{\@ifnextchar[{\olref@i}{\olref@final}}
\def\olref@i[#1]{\@ifnextchar[{\olref@ii}{\olref@final}}
\def\olref@ii[#1]{\@ifnextchar[{\olref@iii}{\olref@final}}
\def\olref@iii[#1]{\olref@final}
\def\olref@final#1{\ref{#1}}
\makeatother
\newcommand*{\olnewpage}{\clearpage}

% ===================== OLP compatibility layer (1) =====================
% Compatibility layer for Open Logic Project sources (pinned commit 1e960beff9ed7835bf3e3f1335e21af3439cd107).
%
% OLP's own style stack (sty/open-logic.sty) needs ntheorem, prooftree, extarrows, mathdots,
% isomath and subfiles; the project's local and host TeX Live installs lack several of those.
% Rather than vendor packages, this file supplies drop-in equivalents for every OLP construct
% that the frozen units actually use. It only rearranges the author's own tokens; no
% mathematical content is added, and nothing here is claimed to be OLP's own file.

\makeatletter

% ---------------- proof trees (OLP's prooftree is a bussproofs derivative) ----------------
% These commands CONSUME their argument and typeset nothing: the bracket labels of OLP's proof
% trees carry rule names. Printing them inline dropped those labels into math mode and produced
% 63 "Missing $" / "allowed only in math mode" errors; swallowing them keeps the brace balance
% and the mode correct. The derivation text between the commands is still printed unchanged.
\providecommand{\AxiomC}[1]{}
\providecommand{\UnaryInfC}[1]{}
\providecommand{\BinaryInfC}[1]{}
\providecommand{\TrinaryInfC}[1]{}
\providecommand{\QuaternaryInfC}[1]{}
\providecommand{\QuinaryInfC}[1]{}
\providecommand{\RightLabel}[1]{}
\providecommand{\LeftLabel}[1]{}
\providecommand{\doubleLine}{}
\providecommand{\noLine}{}
\providecommand{\DisplayProof}{}
\providecommand{\fCenter}{}
\newenvironment{prooftree}{}{}
\newenvironment{ltprooftree}{}{}
\newenvironment{oltableau}{}{}

% ---------------- logical notation ----------------
\providecommand{\lif}{\mathbin{\rightarrow}}
\providecommand{\lforall}{\mathord{\forall}}
\providecommand{\lexists}{\mathord{\exists}}
\providecommand{\ltrue}{\top}
\providecommand{\lfalse}{\bot}
\providecommand{\Ltrue}{\top}
\providecommand{\Lfalse}{\bot}
\providecommand{\Sequent}{\mathrel{\Rightarrow}}
\providecommand{\Proves}{\mathrel{\vdash}}
\providecommand{\Entails}{\mathrel{\models}}
\providecommand{\Subst}[3]{#1[#2/#3]}
\providecommand{\sFmla}[2]{#1(#2)}
\providecommand{\pFmla}[2]{#1(#2)}
\providecommand{\formula}{\relax}
\providecommand{\True}{\top}
\providecommand{\False}{\bot}
\providecommand{\Args}[1]{#1}
\providecommand{\Ts}{\ensuremath{\top}}
\providecommand{\Fi}{\ensuremath{\bot}}
\providecommand{\lfalsum}{\bot}

% ---------------- objects, structures, metatheory ----------------
\providecommand{\Obj}{}
\providecommand{\Struct}{\ensuremath{\mathfrak{A}}}
\providecommand{\Domain}{\ensuremath{|\mathfrak{A}|}}
\providecommand{\Assign}{\ensuremath{s}}
\providecommand{\Sat}[2]{\ensuremath{#1 \models #2}}
\providecommand{\mSat}[3]{\ensuremath{#1 \models #2[#3]}}
\providecommand{\mModel}[2]{\ensuremath{#1 \models #2}}
\providecommand{\Lang}{\ensuremath{\mathcal{L}}}
\providecommand{\Th}{\ensuremath{\mathrm{Th}}}
\providecommand{\Log}{\ensuremath{\mathrm{L}}}
\providecommand{\tuple}[1]{\langle #1 \rangle}
\providecommand{\num}[1]{\underline{#1}}
\providecommand{\cfind}[1]{\varphi_{#1}}
\providecommand{\Nat}{\ensuremath{\mathbb{N}}}
\providecommand{\Mult}{\ensuremath{\mathrm{Mult}}}
\providecommand{\Add}{\ensuremath{\mathrm{Add}}}
\providecommand{\Exp}{\ensuremath{\mathrm{Exp}}}
\providecommand{\Zero}{\ensuremath{0}}
\providecommand{\Succ}{\ensuremath{\mathrm{S}}}
\providecommand{\lambd}{\ensuremath{\lambda}}
\providecommand{\bredpar}{\mathrel{\rightarrow_\beta}}
\providecommand{\bcd}[1]{\ensuremath{\ulcorner #1 \urcorner}}
\providecommand{\Set}[1]{\ensuremath{\{#1\}}}
\providecommand{\indcase}[1]{\par\noindent\emph{Case #1}.\quad}
\providecommand{\indfrm}[1]{\par\noindent\emph{Form #1}.\quad}
\providecommand{\iftag}[1]{#1}
\providecommand{\tagitem}[1]{\item[#1]}
% these five are used with arguments by OLP's proof trees -> swallow the argument
\providecommand{\Elim}[1]{}
\providecommand{\Intro}[1]{}
\providecommand{\LeftR}[1]{}
\providecommand{\RightR}[1]{}
\providecommand{\TAss}[1]{}
\providecommand{\TRule}[1]{}

\makeatother
\providecommand{\olpath}{.}

% ===================== OLP compatibility layer (2) =====================
% --- second compatibility block: notation and shorthands used by the frozen OLP units ---
% Everything here is \providecommand, so anything already defined by LaTeX or the packages wins.
\providecommand{\CutCS}{\mathrm{Cut}}
\providecommand{\LeftSubproofLabel}{\relax}
\providecommand{\RightSubproofLabel}{\relax}
\providecommand{\maeh}{\relax}
\providecommand{\cheight}{\relax}
\providecommand{\ande}{\mathrm{and}_E}
\providecommand{\andi}{\mathrm{and}_I}
\providecommand{\ore}{\mathrm{or}_E}
\providecommand{\ori}{\mathrm{or}_I}
\providecommand{\typeof}{\mathord{:}}
\providecommand{\pValue}[1]{#1}
\providecommand{\Sent}{\mathrm{Sent}}
\providecommand{\cardplus}{\mathbin{+}}
\providecommand{\cardtimes}{\mathbin{\cdot}}
\providecommand{\cardexpo}{\mathbin{\hat{\ }}}
\providecommand{\cardnless}{\not<}
\providecommand{\cardle}{\leq}
\providecommand{\funfromto}[3]{#1\colon #2 \to #3}
\providecommand{\funimage}[1]{\operatorname{im}(#1)}
\providecommand{\fdefined}{\mathord{\downarrow}}
\providecommand{\canonord}{\operatorname{ord}}
\providecommand{\ordtimes}{\mathbin{\cdot}}
\providecommand{\ordeq}[2]{#1 = #2}
\providecommand{\ordneq}[2]{#1 \neq #2}
\providecommand{\ordtype}[1]{\operatorname{ord}(#1)}
\providecommand{\ordsucc}{\mathord{+1}}
\providecommand{\ordunion}{\cup}
\providecommand{\onesphere}{S^1}
\providecommand{\rotationsgroup}{\mathrm{SO}(2)}
\providecommand{\equivrep}[1]{[#1]}
\providecommand{\rlexless}{\prec}
\providecommand{\ZFC}{\mathrm{ZFC}}
\providecommand{\ZF}{\mathrm{ZF}}
\providecommand{\Rat}{\mathbb{Q}}
\providecommand{\Real}{\mathbb{R}}
\providecommand{\PosInt}{\mathbb{Z}^{+}}
\providecommand{\Id}{\mathrm{id}}
\providecommand{\Bin}{\mathrm{Bin}}
\providecommand{\pto}{\rightharpoonup}
\providecommand{\dom}{\operatorname{dom}}
\providecommand{\ran}{\operatorname{ran}}
\providecommand{\len}{\operatorname{len}}
\providecommand{\concat}{\mathbin{{}^\frown}}
\providecommand{\card}[1]{\lvert #1 \rvert}
\providecommand{\Pow}{\mathcal{P}}
\providecommand{\Setabs}[2]{\{\, #1 \mid #2 \,\}}
\providecommand{\comp}{\circ}
\providecommand{\Char}{\chi}
\providecommand{\Prf}{\operatorname{Prf}}
\providecommand{\Prov}{\operatorname{Prov}}
\providecommand{\Proj}{\operatorname{proj}}
\providecommand{\Gn}{\operatorname{gn}}
\providecommand{\gn}{\operatorname{gn}}
\providecommand{\Value}[1]{\operatorname{val}(#1)}
\providecommand{\Deduce}{\mathrm{D}}
\providecommand{\Dual}{\mathrm{d}}
\providecommand{\Ax}{\mathrm{Ax}}
\providecommand{\pSat}[3]{#1 \models #2\,[#3]}
\providecommand{\pAssign}{\sigma}
\providecommand{\bforall}{\forall}
\providecommand{\bexists}{\exists}
\providecommand{\liff}{\mathbin{\leftrightarrow}}
\providecommand{\FV}{\operatorname{FV}}
\providecommand{\UnaryInf}{\ensuremath{\mid}}
\providecommand{\BinaryInf}{\ensuremath{\mid}}
\providecommand{\depth}{\operatorname{depth}}
\providecommand{\oliflabeldef}[3]{#1}
\providecommand{\TMblank}{\sqcup}
\providecommand{\TMstroke}{\mid}
\providecommand{\TMstay}{\mathrm{S}}
\providecommand{\TMright}{\mathrm{R}}
\providecommand{\TMleft}{\mathrm{L}}
\providecommand{\TMendtape}{\triangleright}
\providecommand{\umin}{\mu}
\providecommand{\Sum}{\sum}
% bibliography helpers (natbib): the frozen units cite the shared OLP bibliography
\providecommand{\citep}[1]{\cite{#1}}
\providecommand{\citet}[1]{\cite{#1}}
\providecommand{\citeauthor}[1]{\cite{#1}}
\providecommand{\citeyear}[1]{\cite{#1}}

% ===================== OLP compatibility layer (3) =====================
% --- third compatibility block: remaining notation seen in the frozen OLP units ---
\providecommand{\iso}{\cong}
\providecommand{\nszero}{\mathrm{0}}
\providecommand{\nssucc}{\mathrm{S}}
\providecommand{\nsplus}{\mathbin{+}}
\providecommand{\nssless}{<}
\providecommand{\nsless}{<}
\providecommand{\PL}{\mathrm{PL}}
\providecommand{\Nec}{\mathord{\Box}}
\providecommand{\RK}{\mathrm{RK}}
\providecommand{\ST}{\mathrm{ST}}
\providecommand{\SSubst}[3]{#1[#2/#3]}
\providecommand{\subst}[3]{#1[#2/#3]}
\providecommand{\tagfalse}{\mathrm{F}}
\providecommand{\cutr}{\mathrm{Cut}}
\providecommand{\cardneq}{\neq}
\providecommand{\issucc}[1]{\operatorname{succ}(#1)}
\providecommand{\tup}[1]{\langle #1 \rangle}
\providecommand{\Var}{\mathrm{Var}}
\providecommand{\Bew}{\operatorname{Bew}}
\providecommand{\Con}{\operatorname{Con}}
\providecommand{\Cons}{\operatorname{Cons}}
\providecommand{\pDeduce}{\mathrm{D}}
\providecommand{\psat}{\models}
\providecommand{\den}[1]{\llbracket #1 \rrbracket}
\providecommand{\Sem}[2]{\llbracket #1 \rrbracket_{#2}}
\providecommand{\DomainOf}[1]{\lvert #1 \rvert}
\providecommand{\Form}{\mathrm{Form}}
\providecommand{\Terms}{\mathrm{Term}}
\providecommand{\Atom}{\mathrm{Atom}}
\providecommand{\FVars}[1]{\operatorname{FV}(#1)}
\providecommand{\fv}[1]{\operatorname{fv}(#1)}
\providecommand{\tuplep}[1]{\langle #1 \rangle}
\providecommand{\restrict}[2]{#1 \upharpoonright #2}
\providecommand{\sto}{\rightarrow}
\providecommand{\FmlaSet}[1]{\mathrm{Fmla}_{#1}}
\providecommand{\Thm}{\mathrm{Thm}}
\providecommand{\Eval}{\operatorname{val}}
\providecommand{\blank}{\sqcup}
\providecommand{\blankTape}{\sqcup}
\providecommand{\Consis}{\operatorname{Con}}
\providecommand{\drs}{\mathrm{d}}
\providecommand{\VRule}{\mid}
\providecommand{\Bet}{\mathrm{B}}
\providecommand{\Seg}[1]{\mathrm{Seg}_{#1}}
\providecommand{\Enums}{\mathrm{En}}
\providecommand{\Z}{\mathbb{Z}}
\providecommand{\N}{\mathbb{N}}
\providecommand{\abs}[1]{\lvert #1 \rvert}
\providecommand{\floor}[1]{\lfloor #1 \rfloor}
\providecommand{\ceil}[1]{\lceil #1 \rceil}

% ===================== OLP compatibility layer (4) =====================
% --- fourth compatibility block: harvested from the per-candidate compile logs ---
\providecommand{\beredpar}{\mathrel{\rightarrow_\beta}}
% \bered / \bred are used as infix relations AND with a bracketed annotation, so they need the
% optional-argument chain; \providecommand would leave an argument-less definition in place.
\makeatletter
\def\bered{\@ifnextchar[{\bered@i}{\mathrel{\rightarrow_\beta}}}
\def\bered@i[#1]{\mathrel{\rightarrow_\beta}}
\def\bred{\@ifnextchar[{\bred@i}{\mathrel{\rightarrow_\beta}}}
\def\bred@i[#1]{\mathrel{\rightarrow_\beta}}
\makeatother
\providecommand{\beredone}{\mathrel{=_\beta}}
\providecommand{\bredone}{\mathrel{=_\beta}}
\providecommand{\xredone}{\mathrel{=_\beta}}
\providecommand{\aconvone}{\mathrel{=_\alpha}}
\providecommand{\cardless}{<}
\providecommand{\tagrefs}[1]{#1}
\providecommand{\ordplus}{\mathbin{+}}
\providecommand{\funrestrictionto}[2]{#1 \upharpoonright #2}
\providecommand{\fundefined}{\mathord{\uparrow}}
\providecommand{\Trm}{\mathrm{Trm}}
\providecommand{\OProv}{\operatorname{Prov}}
\providecommand{\OPrf}{\operatorname{Prf}}
\providecommand{\aeq}{\mathrel{=}}
\providecommand{\mClass}[1]{[#1]}
\providecommand{\equivclass}[1]{[#1]}
\providecommand{\pheight}{\operatorname{ht}}
\providecommand{\phy}{\operatorname{ht}}
\providecommand{\Cut}{\mathrm{Cut}}
\providecommand{\QuantRank}{\operatorname{qr}}
\providecommand{\elemequiv}{\mathrel{\equiv}}
\providecommand{\PIso}{\mathrm{PIso}}
\providecommand{\Theory}{\mathrm{Th}}
\providecommand{\supstrict}{\subsetneq}
\providecommand{\trcl}{\operatorname{trcl}}
\providecommand{\cref}[1]{\ref{#1}}
\providecommand{\Discharge}{\mathrm{Dis}}
\providecommand{\Weakening}{\mathrm{Wk}}
\providecommand{\ZFminus}{\mathrm{ZF}^{-}}
\providecommand{\RK}{\mathrm{RK}}
% \Dual is printed both in math mode and (with a bracketed annotation) in running text;
% the math-mode test keeps \mathrm legal in both cases.
\makeatletter
\renewcommand{\Dual}{\@ifnextchar[{\Dual@i}{\ifmmode\mathrm{d}\else$\mathrm{d}$\fi}}
\def\Dual@i[#1]{\ifmmode\mathrm{d}[#1]\else$\mathrm{d}[#1]$\fi}
\makeatother
\providecommand{\RightR}{\mathrm{R}}
\providecommand{\LeftR}{\mathrm{L}}
\providecommand{\Intro}{\mathrm{I}}
\providecommand{\Axiom}{\mathrm{Ax}}
\providecommand{\formula}[1]{#1}
\providecommand{\stared}{\ast}
\providecommand{\oset}[2]{\{\, #1 \mid #2 \,\}}
\providecommand{\setof}[2]{\{\, #1 \mid #2 \,\}}
\providecommand{\pow}[1]{\mathcal{P}(#1)}
\providecommand{\ordpair}[2]{\langle #1, #2 \rangle}
\providecommand{\ordtriple}[3]{\langle #1, #2, #3 \rangle}
\providecommand{\union}{\cup}
\providecommand{\inter}{\cap}
\providecommand{\seq}[1]{\langle #1 \rangle}
\providecommand{\SuchThat}{\mid}
\providecommand{\where}{\mid}
\providecommand{\TVar}{\mathrm{TVar}}
\providecommand{\Bool}{\mathrm{Bool}}
\providecommand{\Val}{\operatorname{Val}}
\providecommand{\FV}{\operatorname{FV}}
\providecommand{\TL}[1]{\mathrm{TL}}
\providecommand{\xstared}{\ast}

% ===================== OLP compatibility layer (5) =====================
% --- fifth compatibility block: rescues the quarantined OLP candidates ---
% The proof-tree and mode fixes live in block 1 (where those macros are first defined); this block
% only adds notation that the quarantined spans still needed and that no earlier block covered.
\providecommand{\ident}{\operatorname{id}}
\providecommand{\setrank}{\operatorname{rank}}
\providecommand{\cardfont}{\mathfrak{c}}
\providecommand{\Complement}{\overline}
\providecommand{\QR}{\mathrm{QR}}
\providecommand{\becd}{\ensuremath{\ulcorner}}
\providecommand{\disjointsum}{\mathbin{\oplus}}
\providecommand{\cardeq}[2]{#1 = #2}
\providecommand{\xred}{\mathrel{\rightarrow}}
\providecommand{\cardless}{<}
\providecommand{\cardplus}{\mathbin{+}}
\providecommand{\cardtimes}{\mathbin{\cdot}}
\providecommand{\cardexpo}{\mathbin{\hat{\ }}}
\providecommand{\cardneq}{\neq}
\providecommand{\cardle}{\leq}
\providecommand{\cardnless}{\not<}
\begin{document}

% ==== OLP file: content/normal-modal-logic/filtrations/S5-decidable.tex ====
% Part: normal-modal-logic
% Chapter: filtrations
% Section: S5-decidable



\olfileid{nml}{fil}{dec}

\olsection{\Log{S5} is Decidable}
The finite model property gives us an easy way to show that systems of
modal logic given by schemas are \emph{decidable} (i.e., that there is
a computable procedure to determine whether !!a{formula} is !!{derivable} in
the system or not).

\begin{thm}
  \Log{S5} is decidable.
\end{thm}

\begin{proof}
  Let $!A$ be given, and suppose the !!{propositional variable}s
  occurring in $!A$ are among $p_1$, \dots, $p_k$.  Since for each
  $n$ there are only finitely many models with $n$ worlds assigning a
  value to $p_1$, \dots, $p_k$, we can enumerate, \emph{in parallel}, all
  the theorems of \Log{S5} by generating proofs in some systematic
  way; and all the models containing $1$, $2$, \dots worlds and checking
  whether $!A$ fails at a world in some such model. Eventually
  one of the two parallel processes will give an answer, as by
  \olref[com][fra]{thm:generaldet} and \olref[fmp]{cor:S5fmp}, either
  $!A$ is !!{derivable} or it fails in a finite universal model.
\end{proof}

The above proof works for \Log{S5} because filtrations of universal
models are automatically universal. The same holds for reflexivity and
seriality, but more work is needed for other properties.


% ==== OLP file: content/normal-modal-logic/filtrations/S5-fmp.tex ====
% Part: normal-modal-logic
% Chapter: filtrations
% Section: S5-fmp



\olfileid{nml}{fil}{fmp}

\olsection{\Log{K} and \Log{S5} have the Finite Model Property}

\begin{defn}
  A system $\Sigma$ of modal logic is said to have the \emph{finite
    model property} if whenever !!a{formula}~$!A$ is true at a world
  in a model of $\Sigma$ then $!A$ is true at a world in a
  \emph{finite} model of~$\Sigma$.
\end{defn}

\begin{prop}\ollabel{prop:K-fmp}
  \Log{K} has the finite model property.
\end{prop}

\begin{proof}
  \Log{K} is the set of valid !!{formula}s, i.e., any model is a model
  of~\Log{K}.  By \olref[fil]{thm:filtrations}, if
  $\mSat{M}{!A}[w]$, then $\mSat{M^*}{!A}[w]$ for any filtration
  of~$\mModel{M}$ through the set $\Gamma$ of sub-!!{formula}s
  of~$!A$. Any !!{formula} only has finitely many sub-!!{formula}s, so
  $\Gamma$ is finite. By \olref[fin]{prop:filt-are-finite},
  $\card{W^*} \le 2^n$, where $n$ is the number of !!{formula}s
  in~$\Gamma$. And since \Log{K} imposes no restriction on models,
  $\mModel{M^*}$ is a \Log{K}-model.
\end{proof}

To show that a logic~\Log{L} has the finite model property via
filtrations it is essential that the filtration of an \Log{L}-model is
itself a \Log{L}-model. Often this requires a fair bit of work, and
not any filtration yields a \Log{L}-model. However, for universal
models, this still holds.

\begin{prop}\ollabel{prop:univ-fin}
  Let $\mClass{U}$ be the class of universal models (see
  \olref[frd][es5]{prop:S5=univ}) and $\mClass{U}_\mathrm{Fin}$ the
  class of all finite universal models. Then any !!{formula}~$!A$ is
  valid in $\mClass{U}$ if and only if it is valid in
  $\mClass{U}_\mathrm{Fin}$.
\end{prop}

\begin{proof}
  Finite universal models are universal models, so the left-to-right
  direction is trivial. For the right-to left direction, suppose
  that~$!A$ is false at some world $w$ in a universal model
  $\mModel{M}$. Let $\Gamma$ contain $!A$ as well as all of its
  subformulas; clearly $\Gamma$ is finite. Take a filtration
  $\mModel{M^*}$ of $\mModel{M}$; then $\mModel{M^*}$ is finite by
  \olref[fin]{prop:filt-are-finite}, and by
  \olref[fil]{thm:filtrations}, $!A$ is false at $[w]$ in
  $\mModel{M^*}$. It remains to observe that $\mModel{M^*}$ is also
  universal: given $u$ and $v$, by hypothesis $Ruv$ and by
  \olref[fil]{defn:filtration}\olref[fil]{defn:filtration-R},
  also $R^*[u][v]$.
\end{proof}

\begin{cor}\ollabel{cor:S5fmp}
  \Log{S5} has the finite model property.
\end{cor}

\begin{proof}
  By \olref[frd][es5]{prop:S5=univ}, if $!A$ is true at a world in
  some reflexive and euclidean model then it is true at a world in a
  universal model. By \olref{prop:univ-fin}, it is true at a world in
  a finite universal model (namely the filtration of the model through
  the set of sub-!!{formula}s of~$!A$). Every universal model is also
  reflexive and euclidean; so $!A$ is true at a world in a finite
  reflexive euclidean model.
\end{proof}

\begin{prob}
  Show that any filtration of a serial or reflexive model is also
  serial or reflexive (respectively).
\end{prob}

\begin{prob}
  Find a non-symmetric (non-transitive, non-euclidean) filtration of a
  symmetric (transitive, euclidean) model.
\end{prob}


% ==== OLP file: content/normal-modal-logic/filtrations/euclidean-filtrations.tex ====
% Part: normal-modal-logic
% Chapter: filtrations
% Section: euclidean-filtrations



\olfileid{nml}{fil}{euc}

\olsection{Filtrations of Euclidean Models}

The approach of \olref[acc]{sec} does not work in the case of models
that are euclidean or serial and euclidean. Consider the model at the
top of \olref{fig:ser-eucl}, which is both euclidean and serial. Let
$\Gamma = \{p, \Box p \}$. When taking a filtration through $\Gamma$,
then $[w_1] = [w_3]$ since $w_1$ and $w_3$ are the only worlds that
agree on~$\Gamma$. Any filtration will also have the arrow inherited
from $\mModel{M}$, as depicted in \olref{fig:ser-eucl2}. That model
isn't euclidean. Moreover, we cannot add arrows to that model in order
to make it euclidean. We would have to add double arrows between
$[w_2]$ and $[w_4]$, and then also between $w_2$ and~$w_5$. But $\Box
p$ is supposed to be true at~$w_2$, while $p$ is false at~$w_5$.

\begin{figure}[htpb]
  \centering
  \begin{tikzpicture}[modal]
    \node[world] (w1)
          [label={left: \mFalse{p}},
            label={below: $\mSat{{}}{\Box p}$}] {$w_1$}; 
    \node[world] (w2)
          [label={right: \mTrue{p}},
            label = {below: $\mSat{{}}{\Box p}$},
            right=of w1] {$w_2$}; 
    \draw[->] (w1) to (w2);
    \draw[reflexive above] (w2) to (w2);  
    \node[world] (w3)
          [label={left: \mFalse{p}},
            label={below: $\mSat{{}}{\Box p}$},
            below=of w1] {$w_3$}; 
    \node[world] (w4)
          [label={right:\mTrue{p}},
            label={below: $\mSat/{{}}{\Box p}$},
            right=of w3] {$w_4$};
    \draw[->] (w3) to (w4);
    \draw[reflexive above] (w4) to (w4);     
    \node[world] (w5)
          [label={right: \mFalse{p}},
            label=below:$\mSat/{{}}{\Box p}$,
            right=of w4] {$w_5$}; 
    \draw[->, bend left] (w4) to (w5);
    \draw[->, bend left] (w5) to (w4);
    \draw[reflexive above] (w5) to (w5);     
  \end{tikzpicture}
  \caption{A serial and euclidean model.}
  \ollabel{fig:ser-eucl}
\end{figure}

\begin{figure}[ht]
  \centering
  \begin{tikzpicture}[modal,every text node part/.style={align=center}]
    \node[world] (w1)
          [label={left: \mFalse{p}},
            label={right:$[w_1]=[w_3]$},
            label={below: $\mSat{{}}{\Box p}$}] {$[w_1]$}; 
    \node[world] (w2)
         [label={right: \mTrue{p}},
           label = {below: $\mSat{{}}{\Box p}$},
           right=of w1, yshift=1.5cm] {$[w_2]$}; 
    \draw[->] (w1) to (w2);
    \node[world] (w4)
         [label={right:\mTrue{p}},
           label={below: $\mSat/{{}}{\Box p}$},
           right=of w1,yshift=-1.5cm] {$[w_4]$};
    \draw[->] (w1) to (w4);
    \draw[reflexive above] (w4) to (w4);     
    \node[world] (w5)
         [label={right: \mFalse{p}},
             label=below:$\mSat/{{}}{\Box p}$,
             right=of w4] {$[w_5]$}; 
    \draw[->, bend left] (w4) to (w5);
    \draw[->, bend left] (w5) to (w4);
    \draw[reflexive above] (w5) to (w5);     
  \end{tikzpicture}
  \caption{The filtration of the model in \olref{fig:ser-eucl}.}
  \ollabel{fig:ser-eucl2}
\end{figure}

In particular, to obtain a euclidean filtration it is not enough to
consider filtrations through arbitrary $\Gamma$'s closed under
sub-!!{formula}s. Instead we need to consider sets $\Gamma$ that are
\emph{modally closed} (see \olref[pre]{defn:modallyclosed}). Such sets
of sentences are infinite, and therefore do not immediately yield a
finite model property or the decidability of the corresponding system.

\begin{thm}\ollabel{thm:modal-closed-filt}
  Let $\Gamma$ be modally closed, $\mModel{M}=\tuple{W,R,V}$, and
  $\mModel{M^*} = \tuple{W^*,R^*,V^*}$ be a coarsest filtration of
  $\mModel{M}$.
  \begin{enumerate}
  \item If $\mModel{M}$ is symmetric, so is $\mModel{M^*}$.
  \item If $\mModel{M}$ is transitive, so is $\mModel{M^*}$.
  \item If $\mModel{M}$ is euclidean, so is $\mModel{M^*}$.
  \end{enumerate}
\end{thm}

\begin{proof}
\begin{enumerate}
  \item Exercise. Use the fact that \Ax{B} and $\Ax{B_\Diamond}$ are
    valid in all symmetric models.
  \item If $\mModel{M^*}$ is a coarsest filtration, then by definition
    $R^*[u][v]$ holds if and only if $C_1(u,v)$. For transitivity,
    suppose $C_1(u,v)$ and $C_1(v,w)$; we have to show
    $C_1(u,w)$. \iftag{prvBox}{Suppose $\mSat{M}{\Box !A}[u]$; then
      $\mSat{M}{\Box\Box!A}[u]$ since \Ax{4} is valid in all
      transitive models; since $\Box\Box!A \in \Gamma$ by closure,
      also by $C_1(u,v)$, $\mSat{M}{\Box!A}[v]$ and by $C_1(v,w)$,
      also $\mSat{M}{!A}[w]$. }{}\iftag{prvDiamond}{Suppose
      $\mSat{M}{!A}[w]$; then $\mSat{M}{\Diamond !A}[v]$ by
      $C_1(v,w)$, since $\Diamond !A \in \Gamma$ by modal closure. By
      $C_1(u,v)$, we get $\mSat{M}{\Diamond\Diamond !A}[u]$ since
      $\Diamond\Diamond!A \in \Gamma$ by modal closure. Since
      $\Ax{4}_\Diamond$ is valid in all transitive models,
      $\mSat{M}{\Diamond!A}[u]$.}{}
  \item Exercise. Use the fact that both \Ax{5} and $\Ax{5_\Diamond}$
    are valid in all euclidean models.
\end{enumerate}
\end{proof}

\begin{prob}
  Complete the proof of \olref[nml][fil][euc]{thm:modal-closed-filt}.
\end{prob}


% ==== OLP file: content/normal-modal-logic/filtrations/examples-of-filtrations.tex ====
% Part: normal-modal-logic
% Chapter: filtrations
% Section: examples-of-filtrations



\olfileid{nml}{fil}{exf}

\olsection{Examples of Filtrations}

We have not yet shown that there are any filtrations. But indeed, for
any model $\mModel{M}$, there are many filtrations of $\mModel{M}$
through $\Gamma$. We identify two, in particular: the finest and
coarsest filtrations. Filtrations of the same models will differ in
their accessibility relation (as \olref[fil]{defn:filtration} stipulates
directly what $W^*$ and~$V^*$ should be). The finest filtration
will have as few related worlds as possible, whereas the coarsest will
have as many as possible.

\begin{defn}
  Where $\Gamma$ is closed under subformulas, the \emph{finest}
  filtration $\mModel{M^*}$ of a model $\mModel{M}$ is defined by putting:
  \[
  R^*[u][v] \quad \text{if and only if} \quad \exists u'\in [u] \;
  \exists v' \in [v] : Ru'v'.
  \]
\end{defn}

\begin{prop}\ollabel{prop:finest}
  The finest filtration $\mModel{M^*}$ is indeed a filtration. 
\end{prop}

\begin{proof}
  We need to check that $R^*$, so defined, satisfies
  \olref[fil]{defn:filtration}\olref[fil]{defn:filtration-R}. We check
  the three conditions in turn.

  If $Ruv$ then since $u \in [u]$ and $v \in [v]$, also $R^*[u][v]$,
  so \olref[fil]{defn:filtration-R1} is satisfied.
  
  \iftag{prvBox}{\iftag{probBox}{We leave the verification of
      \olref[fil]{defn:filtration-R2} as an exercise.}{For
        \olref[fil]{defn:filtration-R2}, suppose $\Box!A \in \Gamma$,
        $R^*[u][v]$, and $\mSat{M}{\Box!A}[u]$. By definition of
        $R^*$, there are $u' \equiv u$ and $v' \equiv v$ such that
        $Ru'v'$. Since $u$ and $u'$ agree on $\Gamma$, also
        $\mSat{M}{\Box!A}[u']$, so that $\mSat{M}{!A}[v']$. By closure
        of $\Gamma$ under sub-!!{formula}s, $v$ and $v'$ agree on
        $!A$, so $\mSat{M}{!A}[v]$, as desired.}}{}
  
  \iftag{prvDiamond}{\iftag{probDiamond}{We leave the verification of
      \olref[fil]{defn:filtration-R3} as an exercise.}{To verify
        \olref[fil]{defn:filtration-R3}, suppose $\Diamond!A \in
        \Gamma$, $R^*[u][v]$, and $\mSat{M}{!A}[v]$.  By definition of
        $R^*$, there are $u' \equiv u$ and $v' \equiv v$ such that
        $Ru'v'$. Since $v$ and $v'$ agree on~$\Gamma$, and $\Gamma$ is
        closed under sub-!!{formula}s, also $\mSat{M}{!A}[v']$, so
        that $\mSat{M}{\Diamond !A}[u']$.  Since $u$ and $u'$ also
        agree on $\Gamma$, $\mSat{M}{\Diamond !A}[u]$.}}{}
\end{proof}

\begin{probtag}{probBox,probDiamond}
  Complete the proof of \olref[nml][fil][exf]{prop:finest}.
\end{probtag}

\begin{defn}
  Where $\Gamma$ is closed under subformulas, the \emph{coarsest}
  filtration~$\mModel{M^*}$ of a model~$\mModel{M}$ is defined by
  putting $R^*[u][v]$ if and only if
  \iftag{notprvBox,notprvDiamond}{the following condition is
    met:}{\emph{both} of the following conditions are met:}
  \begin{tagenumerate}{prvBox,prvDiamond}
  \tagitem{prvBox}{\ollabel{defn:coarsest-Box}If $\Box!A \in \Gamma$
    and $\mSat{M}{\Box!A}[u]$ then
    $\mSat{M}{!A}[v]$\iftag{prvDiamond}{;}{.}}{}
  \tagitem{prvDiamond}{\ollabel{defn:coarsest-Diamond}If $\Diamond!A
    \in \Gamma$ and $\mSat{M}{!A}[v]$ then
    $\mSat{M}{\Diamond!A}[u]$.}{}
  \end{tagenumerate}
\end{defn}

\begin{prop}
  The coarsest filtration $\mModel{M^*}$ is indeed a filtration. 
\end{prop}

\begin{proof}
  Given the definition of $R^*$, the only condition that is left to
  verify is the implication from $Ruv$ to $R^*[u][v]$. So assume
  $Ruv$. \iftag{prvBox}{Suppose $\Box!A \in \Gamma$ and
    $\mSat{M}{\Box!A}[u]$; then obviously
    $\mSat{M}{!A}[v]$\iftag{prvDiamond}{, and
      \olref{defn:coarsest-Box} is
      satisfied. }{.}}{}\iftag{prvDiamond}{Suppose $\Diamond!A \in
    \Gamma$ and $\mSat{M}{!A}[v]$. Then $\mSat{M}{\Diamond!A}[u]$
    since~$Ruv$\iftag{prvBox}{, and \olref{defn:coarsest-Diamond} is
      satisfied}{}.}{}
\end{proof}

\begin{ex}
  Let $W = \PosInt$, $Rnm$ iff $m = n + 1$, and $V(p) = \Setabs{2n}{n \in
    \Nat}$. The model $\mModel{M} = \tuple{W, R, V}$ is depicted in
  \olref{fig:ex-filtration}. The worlds are $1$, $2$, etc.; each
  world can access exactly one other world---its successor---and $p$ is
  true at all and only the even numbers.

  \begin{figure}
    \centering
    \begin{tikzpicture}[modal]
      \node[world] (1) [label=below:\mFalse{p}] {$1$};
      \node[world] (2) [label=below:\mTrue{p}, right=of 1] {$2$};
      \node[world] (3) [label=below:\mFalse{p}, right=of 2] {$3$};
      \node[world] (4) [label=below:\mTrue{p}, right=of 3] {$4$};
      \node[phantom] (5) [right=of 4] {};
      \draw[->] (1) to (2);
      \draw[->] (2) to (3);
      \draw[->] (3) to (4);
      \draw[dotted] (4) to (5);
    \end{tikzpicture}

    \begin{tikzpicture}[modal]
      \node[world] (1) [label=below:\mFalse{p}] {$[1]$};
      \node[world] (2) [label=below:\mTrue{p}, right=of 1] {$[2]$};
      \draw[->,bend left] (1) to (2);
      \draw[->,bend left] (2) to (1);
      
      \node[world] (3) [label=below:\mFalse{p},right=of 2] {$[1]$};
      \node[world] (4) [label=below:\mTrue{p}, right=of 3] {$[2]$};
      \draw[->,bend left] (3) to (4);
      \draw[->,bend left] (4) to (3);
      \draw[->,reflexive above] (4) to (4);
    \end{tikzpicture}
    \caption{An infinite model and its filtrations.}
    \ollabel{fig:ex-filtration}
  \end{figure}
  
  Now let $\Gamma$ be the set of sub-!!{formula}s of~$\Box p \lif p$,
  i.e., $\{p, \Box p, \Box p \lif p\}$. $p$ is true at all and only
  the even numbers, $\Box p$ is true at all and only the odd numbers,
  so $\Box p \lif p$ is true at all and only the even numbers. In
  other words, every odd number makes $\Box p$ true but $p$ and $\Box
  p \lif p$ false; every even number makes $p$ and $\Box p \lif p$
  true, but $\Box p$ false. So $W^* = \{ [1], [2] \}$, where $[1] =
  \{1, 3, 5, \dots\}$ and $[2] = \{2, 4, 6, \dots\}$. Since $2 \in
  V(p)$, $[2] \in V^*(p)$; since $1 \notin V(p)$, $[1] \notin
  V^*(p)$. So $V^*(p) = \{[2]\}$.

  Any filtration based on $W^*$ must have an accessibility relation
  that includes $\tuple{[1], [2]}, \tuple{[2],[1]}$: since $R12$, we
  must have $R^*[1][2]$ by
  \olref[fil]{defn:filtration}\olref[fil]{defn:filtration-R1}, and
  since $R23$ we must have $R^*[2][3]$, and $[3]=[1]$. It cannot
  include $\tuple{[1],[1]}$: if it did, we'd have $R^*[1][1]$,
  $\mSat{M}{\Box p}[1]$ but $\mSat/{M}{p}[1]$, contradicting
  \olref[fil]{defn:filtration-R2}. Nothing requires or rules out that
  $R^*[2][2]$. So, there are two possible filtrations of~$\mModel{M}$,
  corresponding to the two accessibility relations
  \[
  \{\tuple{[1],[2]}, \tuple{[2],[1]}\} \text{ and }
  \{\tuple{[1],[2]}, \tuple{[2],[1]}, \tuple{[2],[2]}\}.
  \]
  In either case, $p$ and $\Box p \lif p$ are false and $\Box p$ is
  true at~$[1]$; $p$ and $\Box p \lif p$ are true and $\Box p$ is
  false at~$[2]$.
\end{ex}


\begin{prob}
  Consider the following model $\mModel{M} = \tuple{W, R, V}$ where $W
  = \Setabs{0\sigma}{\sigma \in \Bin^*}$, the set of sequences of $0$s
  and~$1$s starting with~$0$, with $R\sigma\sigma'$ iff $\sigma' =
  \sigma 0$ or $\sigma' = \sigma 1$, and $V(p) = \Setabs{\sigma
    0}{\sigma \in \Bin^*}$ and $V(q) = \Setabs{\sigma 1}{\sigma \in
    \Bin^* \setminus \{1\}}$. Here's a picture:
  \begin{center}
    \begin{tikzpicture}[modal,
        node distance=2cm
      ]

      \node[world] (0) [label={below:\mTrue{p}\\\mFalse{q}},font=\tiny] {$0$};
      \node[world] (00) [label={below:\mTrue{p}\\\mFalse{q}},
        above right=of 0,font=\tiny] {$00$};
      \node[world] (000) [label={below:\mTrue{p}\\\mFalse{q}},
        right=of 00, yshift=.8cm, font=\tiny] {$000$};
      \node[phantom] (0000) [right=of 000, yshift=.5cm] {};
      \node[phantom] (0001) [right=of 000, yshift=-.5cm] {};
      \node[world] (001) [label={below:\mFalse{p}\\\mTrue{q}},
        right=of 00, yshift=-.8cm, font=\tiny] {$001$};
      \node[phantom](0010) [right=of 001, yshift=.5cm] {};
      \node[phantom](0011) [right=of 001, yshift=-.5cm] {};
      \node[world] (01) [label={below:\mFalse{p}\\\mTrue{q}},
        below right=of 0,font=\tiny] {$01$};
      \node[world] (010) [label={below:\mTrue{p}\\\mFalse{q}},
        right=of 01, yshift=.8cm,font=\tiny] {$010$};
      \node[phantom] (0100) [right=of 010, yshift=.5cm] {};
      \node[phantom] (0101) [right=of 010, yshift=-.5cm] {};
      \node[world] (011) [label={below:\mFalse{p}\\\mTrue{q}},
        right=of 01, yshift=-.8cm,font=\tiny] {$011$};
      \node[phantom] (0110) [right=of 011, yshift=.5cm] {};
      \node[phantom] (0111) [right=of 011, yshift=-.5cm] {};

      \draw[->] (0) to (00);
      \draw[->] (0) to (01);
      \draw[->] (00) to (000);
      \draw[dotted] (000) to (0000);
      \draw[dotted] (000) to (0001);
      \draw[->] (00) to (001);
      \draw[dotted] (001) to (0010);
      \draw[dotted] (001) to (0011);
      \draw[->] (01) to (010);
      \draw[dotted] (010) to (0100);
      \draw[dotted] (010) to (0101);
      \draw[->] (01) to (011);
      \draw[dotted] (011) to (0110);
      \draw[dotted] (011) to (0111);
    \end{tikzpicture}
  \end{center}
  We have $\mSat/{M}{\Box(p \lor q) \lif (\Box p \lor \Box q)}[w]$ for
  every~$w$.

  Let $\Gamma$ be the set of sub-!!{formula}s of~$\Box(p \lor q) \lif
  (\Box p \lor \Box q)$. What are $W^*$ and~$V^*$?  What is the
  accessibility relation of the finest filtration of~$\mModel{M}$? Of the
  coarsest?
\end{prob}


% ==== OLP file: content/normal-modal-logic/filtrations/filtrations-def.tex ====
% Part: normal-modal-logic
% Chapter: filtrations
% Section: filtrations-def



\olfileid{nml}{fil}{fil}

\olsection{Filtrations}

Rather than define ``the'' filtration of $\mModel{M}$ through
$\Gamma$, we define when a model~$\mModel{M^*}$ counts as a filtration
of~$\mModel{M}$. All filtrations have the same set of worlds~$W^*$ and
the same valuation~$V^*$. But different filtrations may have different
accessibility relations~$R^*$. To count as a filtration, $R^*$ has to
satisfy a number of conditions, however. These conditions are exactly
what we'll require to prove the main result, namely that
$\mSat{M}{!A}[w]$ iff $\mSat{M^*}{!A}[{[w]}]$, provided $!A \in
\Gamma$.

\begin{defn}\ollabel{defn:filtration}
  Let $\Gamma$ be closed under subformulas and $\mModel{M} = \tuple{W,
    R, V}$. A \emph{filtration of $\mModel{M}$ through $\Gamma$} is
  any model $\mModel{M^*} = \tuple{W^*,R^*,V^*}$, where:
  \begin{enumerate}
  \item $W^* = \Setabs{[w]}{w \in W}$;
  \item \ollabel{defn:filtration-R}%
    For any $u,v \in W$:
    \begin{enumerate}
    \item \ollabel{defn:filtration-R1}%
      If $Ruv$ then $R^*[u][v]$;
    \tagitem{prvBox}{\ollabel{defn:filtration-R2}%
      If $R^*[u][v]$ then for any $\Box !A \in \Gamma$, if
      $\mSat{M}{\Box!A}[u]$ then $\mSat{M}{!A}[v]$}{};
    \tagitem{prvDiamond}{\ollabel{defn:filtration-R3}%
      If $R^*[u][v]$ then for any $\Diamond !A \in \Gamma$, if
      $\mSat{M}{!A}[v]$ then $\mSat{M}{\Diamond!A}[u]$.}{}
    \end{enumerate}
  \item $V^*(p) = \Setabs{[u]}{u \in V(p)}$.
  \end{enumerate}
\end{defn}

It's worthwhile thinking about what $V^*(p)$ is: the set consisting of
the equivalence classes~$[w]$ of all worlds~$w$ where $p$ is true
in~$\mModel{M}$. On the one hand, if $w \in V(p)$, then $[w] \in
V^*(p)$ by that definition. However, it is not necessarily the case
that if $[w] \in V^*(p)$, then $w \in V(p)$. If $[w] \in V^*(p)$ we
are only guaranteed that $[w] = [u]$ for \emph{some}~$u \in V(p)$. Of
course, $[w] = [u]$ means that $w \equiv u$.  So, when $[w] \in
V^*(p)$ we can (only) conclude that $w \equiv u$ for some $u \in
V(p)$.

\begin{thm}\ollabel{thm:filtrations}
  If $\mModel{M^*}$ is a filtration of $\mModel{M}$ through $\Gamma$,
  then for every $!A \in \Gamma$ and $w \in W$, we have
  $\mSat{M}{!A}[w]$ if and only if $\mSat{M^*}{!A}[{[w]}]$.
\end{thm}

\begin{proof}
  By induction on $!A$, using the fact that $\Gamma$ is closed under
  subformulas. Since $!A \in \Gamma$ and $\Gamma$ is closed under
  sub-!!{formula}s, all sub-!!{formula}s of~$!A$ are also $\in
  \Gamma$. Hence in each inductive step, the induction hypothesis
  applies to the sub-!!{formula}s of~$!A$.
  \begin{enumerate}
    \tagitem{prvFalse}{\indcase{!A}{\lfalse}{Neither
        $\mSat{M}{\indfrm}[w]$ nor $\mSat{M^*}{\indfrm}[{[w]}]$.}}{}
    \tagitem{prvTrue}{\indcase{!A}{\ltrue}{Both
        $\mSat{M}{\indfrm}[w]$ and $\mSat{M^*}{\indfrm}[{[w]}]$.}}{}
    \item \indcase{!A}{p}{The left-to-right direction is immediate, as
      $\mSat{M}{\indfrm}[w]$ only if $w \in V(p)$, which implies $[w]
      \in V^*(p)$, i.e., $\mSat{M^*}{\indfrm}[{[w]}]$. Conversely,
      suppose $\mSat{M^*}{\indfrm}[{[w]}]$, i.e., $[w] \in
      V^*(p)$. Then for some $v \in V(p)$, $w \equiv v$. Of course
      then also $\mSat{M}{p}[v]$. Since $w \equiv v$, $w$ and $v$ make
      the same !!{formula}s from~$\Gamma$ true. Since by assumption $p
      \in \Gamma$ and $\mSat{M}{p}[v]$, $\mSat{M}{\indfrm}[w]$.}

    \tagitem{prvNot}{\iftag{probNot}{Exercise.}{\indcase{!A}{\lnot
          !B}{$\mSat{M}{\indfrm}[w]$ iff $\mSat/{M}{!B}[w]$. By
          induction hypothesis, $\mSat/{M}{!B}[w]$ iff
          $\mSat/{M^*}{!B}[{[w]}]$.  Finally, $\mSat/{M^*}{!B}[{[w]}]$
          iff $\mSat{M^*}{\indfrm}[{[w]}]$.}}}{}
        
    \tagitem{prvAnd}{\iftag{probAnd}{Exercise.}{\indcase{!A}{(!B \land
          !C)}{$\mSat{M}{\indfrm}[w]$ iff $\mSat{M}{!B}[w]$ and
          $\mSat{M}{!C}[w]$. By induction hypothesis, $\mSat{M}{!B}[w]$
          iff $\mSat{M^*}{!B}[{[w]}]$, and $\mSat{M}{!C}[w]$ iff
          $\mSat{M^*}{!C}[{[w]}]$. And $\mSat{M^*}{\indfrm}[{[w]}]$
          iff $\mSat{M^*}{!B}[{[w]}]$ and
          $\mSat{M^*}{!C}[{[w]}]$.}}}{}

    \tagitem{prvOr}{\iftag{probOr}{Exercise.}{\indcase{!A}{(!B \lor
          !C)}{$\mSat{M}{\indfrm}[w]$ iff $\mSat{M}{!B}[w]$ or
          $\mSat{M}{!C}[w]$. By induction hypothesis, $\mSat{M}{!B}[w]$
          iff $\mSat{M^*}{!B}[{[w]}]$, and $\mSat{M}{!C}[w]$ iff
          $\mSat{M^*}{!C}[{[w]}]$. And $\mSat{M^*}{\indfrm}[{[w]}]$
          iff $\mSat{M^*}{!B}[{[w]}]$ or
          $\mSat{M^*}{!C}[{[w]}]$.}}}{}

    \tagitem{prvIf}{\iftag{probIf}{Exercise.}{\indcase{!A}{(!B \lif
          !C)}{$\mSat{M}{\indfrm}[w]$ iff $\mSat/{M}{!B}[w]$ or
          $\mSat{M}{!C}[w]$. By induction hypothesis, $\mSat{M}{!B}[w]$
          iff $\mSat{M^*}{!B}[{[w]}]$, and $\mSat{M}{!C}[w]$ iff
          $\mSat{M^*}{!C}[{[w]}]$. And $\mSat{M^*}{\indfrm}[{[w]}]$
          iff $\mSat/{M^*}{!B}[{[w]}]$ or
          $\mSat{M^*}{!C}[{[w]}]$.}}}{}

    \tagitem{prvIff}{\iftag{probIff}{Exercise.}{\indcase{!A}{(!B \liff
          !C)}{$\mSat{M}{\indfrm}[w]$ iff $\mSat{M}{!B}[w]$ and
          $\mSat{M}{!C}[w]$, or $\mSat/{M}{!B}[w]$ and
          $\mSat/{M}{!C}[w]$. By induction hypothesis, $\mSat{M}{!B}[w]$
          iff $\mSat{M^*}{!B}[{[w]}]$, and $\mSat{M}{!C}[w]$ iff
          $\mSat{M^*}{!C}[{[w]}]$. And $\mSat{M^*}{\indfrm}[{[w]}]$
          iff $\mSat{M^*}{!B}[{[w]}]$ and
          $\mSat{M^*}{!C}[{[w]}]$, or $\mSat/{M^*}{!B}[{[w]}]$ and
          $\mSat/{M^*}{!C}[{[w]}]$.}}}{}

    \tagitem{prvBox}{\iftag{probBox}{Exercise.}{\indcase{!A}{\Box
          !B}{Suppose $\mSat{M}{\indfrm}[w]$; to show that
          $\mSat{M^*}{\indfrm}[{[w]}]$, let $v$ be such that
          $R^*[w][v]$. From
          \olref{defn:filtration}\olref{defn:filtration-R2}, we have
          that $\mSat{M}{!B}[v]$, and by inductive hypothesis
          $\mSat{M^*}{!B}[{[v]}]$. Since $v$ was arbitrary,
          $\mSat{M^*}{\indfrm}[{[w]}]$ follows.

          Conversely, suppose $\mSat{M^*}{\indfrm}[{[w]}]$ and let $v$
          be arbitrary such that $Rwv$. From
          \olref{defn:filtration}\olref{defn:filtration-R1}, we have
          $R^*[w][v]$, so that $\mSat{M^*}{!B}[{[v]}]$; by inductive
          hypothesis $\mSat{M}{!B}[v]$, and since $v$ was arbitrary,
          $\mSat{M}{\indfrm}[w]$.}}}{}
    
    \tagitem{prvDiamond}{\iftag{probDiamond}{Exercise.}{\indcase{!A}{\Diamond
          !B}{Suppose $\mSat{M}{\indfrm}[w]$. Then for some $v \in W$,
          $Rwv$ and $\mSat{M}{!B}[v]$.  By inductive hypothesis
          $\mSat{M^*}{!B}[{[v]}]$, and by
          \olref{defn:filtration}\olref{defn:filtration-R1}, we have
          $R^*[w][v]$. Thus, $\mSat{M^*}{\indfrm}[{[w]}]$.

          Now suppose $\mSat{M^*}{\indfrm}[{[w]}]$. Then for some $[v]
          \in W^*$ with $R^*[w][v]$, $\mSat{M^*}{!B}[{[v]}]$.  By
          inductive hypothesis $\mSat{M}{!B}[v]$.  By
          \olref{defn:filtration}\olref{defn:filtration-R3}, we have
          that $\mSat{M}{\indfrm}[w]$.}}}{}
  \end{enumerate}
\end{proof}

\begin{prob}
  Complete the proof of \olref[nml][fil][fil]{thm:filtrations}
\end{prob}

What holds for truth at worlds in a model also holds for truth in a
model and validity in a class of models.

\begin{cor}
  Let $\Gamma$ be closed under subformulas. Then:
  \begin{enumerate}
  \item If $\mModel{M^*}$ is a filtration of $\mModel{M}$ through
    $\Gamma$ then for any $!A \in \Gamma$: $\mSat{M}{!A}$ if and only
    if $\mSat{M^*}{!A}$.
  \item If $\mClass{C}$ is a class of models and $\Gamma(\mClass{C})$
    is the class of $\Gamma$-filtrations of models in $\mClass{C}$,
    then any !!{formula}~$!A \in \Gamma$ is valid in $\mClass{C}$ if
    and only if it is valid in $\Gamma(\mClass{C})$.
  \end{enumerate}
\end{cor}




% ==== OLP file: content/normal-modal-logic/filtrations/filtrations.tex ====
% Part: modal-logic
% Chapter: filtrations



\olchapter{nml}{fil}{Filtrations and Decidability}

\olimport{introduction}
\olimport{preliminaries}
\olimport{filtrations-def}
\olimport{examples-of-filtrations}
\olimport{finite}
\olimport{S5-fmp}
\olimport{S5-decidable}
\olimport{more-filtrations}
\olimport{euclidean-filtrations}

\OLEndChapterHook



% ==== OLP file: content/normal-modal-logic/filtrations/finite.tex ====
% Part: normal-modal-logic
% Chapter: filtrations
% Section: finite



\olfileid{nml}{fil}{fin}

\olsection{Filtrations are Finite}

We've defined filtrations for any set~$\Gamma$ that is closed under
sub-!!{formula}s. Nothing in the definition itself guarantees that
filtrations are finite. In fact, when $\Gamma$ is infinite (e.g., is
the set of all !!{formula}s), it may well be infinite. However, if
$\Gamma$ is finite (e.g., when it is the set of sub-!!{formula}s of a
given !!{formula}~$!A$), so is any filtration through~$\Gamma$.

\begin{prop}\ollabel{prop:filt-are-finite}
  If $\Gamma$ is finite then any filtration $\mModel{M^*}$ of a model
  $\mModel{M}$ through $\Gamma$ is also finite.
\end{prop}

\begin{proof}
  The size of $W^*$ is the number of different classes~$[w]$ under the
  equivalence relation~$\equiv$. Any two worlds $u$, $v$ in such
  class---that is, any $u$ and $v$ such that~$u \equiv v$---agree on
  all !!{formula}s~$!A$ in~$\Gamma$, $!A \in \Gamma$ either $!A$ is
  true at both $u$ and $v$, or at neither. So each class~$[w]$
  corresponds to subset of~$\Gamma$, namely the set of all $!A \in
  \Gamma$ such that $!A$ is true at the worlds in~$[w]$. No two
  different classes $[u]$ and $[v]$ correspond to the same subset
  of~$\Gamma$. For if the set of !!{formula}s true at $u$ and that of
  !!{formula}s true at $v$ are the same, then $u$ and $v$ agree on all
  formulas in~$\Gamma$, i.e., $u \equiv v$. But then $[u] = [v]$.  So,
  there is !!a{injective} function from $W^*$ to $\Pow{\Gamma}$, and
  hence $\card{W^*} \le \card{\Pow{\Gamma}}$.  Hence if $\Gamma$
  contains $n$ sentences, the cardinality of $W^*$ is no greater
  than~$2^n$.
\end{proof}




% ==== OLP file: content/normal-modal-logic/filtrations/introduction.tex ====
% Part: normal-modal-logic
% Chapter: filtrations
% Section: introduction



\olfileid{nml}{fil}{int}

\olsection{Introduction}

One important question about a logic is always whether it is
decidable, i.e., if there is an effective procedure which will answer
the question ``is this !!{formula} valid.'' Propositional logic is
decidable: we can effectively test if !!a{formula} is a tautology by
constructing a truth table, and for a given !!{formula}, the truth
table is finite.  But we can't obviously test if a modal formula is
true in all models, for there are infinitely many of them.  We can
list all the finite models relevant to a given formula, since only the
assignment of subsets of worlds to !!{propositional variable}s which
actually occur in the !!{formula} are relevant. If the accessibility
relation is fixed, the possible different assignments $V(p)$ are just
all the subsets of~$W$, and if $\card{W} = n$ there are $2^n$ of
those. If our !!{formula}~$!A$ contains $m$ !!{propositional
  variable}s there are then $2^{nm}$ different models with~$n$
worlds. For each one, we can test if~$!A$ is true at all worlds,
simply by computing the truth value of~$!A$ in each. Of course, we
also have to check all possible accessibility relations, but there are
only finitely many relations on~$n$ worlds as well (specifically, the
number of subsets of $W \times W$, i.e., $2^{n^2}$.

If we are not interested in the logic~\Log{K}, but a logic defined by
some class of models (e.g., the reflexive transitive models), we also
have to be able to test if the accessibility relation is of the right
kind. We can do that whenever the frames we are interested in are
definable by modal !!{formula}s (e.g., by testing if \Ax{T} and \Ax{4}
valid in the frame). So, the idea would be to run through all the
finite frames, test each one if it is a frame in the class we're
interested in, then list all the possible models on that frame and
test if $!A$ is true in each. If not, stop: $!A$ is not valid in the
class of models of interest.

There is a problem with this idea: we don't know when, if ever, we can
stop looking. If the !!{formula} has a finite countermodel, our procedure will
find it. But if it has no finite countermodel, we won't get an
answer. The !!{formula} may be valid (no countermodels at all), or it
may have only an infinite countermodel, which we'll never look at. This
problem can be overcome if we can show that every !!{formula} that has
a countermodel has a finite countermodel. If this is the case we say
the logic has the \emph{finite model property}.

But how would we show that a logic has the finite model property? One
way of doing this would be to find a way to turn an infinite
(counter)model of~$!A$ into a finite one. If that can be done, then
whenever there is a model in which $!A$ is not true, then the
resulting finite model also makes $!A$ not true. That finite model
will show up on our list of all finite models, and we will eventually
determine, for every formula that is not valid, that it isn't. Our
procedure won't terminate if the formula is valid. If we can show in
addition that there is some maximum size that the finite model our
procedure provides can have, and that this maximum size depends only
on the !!{formula}~$!A$, we will have a size up to which we have to
test finite models in our search for countermodels. If we haven't
found a countermodel by then, there are none. Then our procedure will,
in fact, decide the question ``is $!A$ valid?'' for any
!!{formula}~$!A$.

A strategy that often works for turning infinite structures into
finite structures is that of ``identifying'' !!{element}s of the
structure which behave the same way in relevant respects. If there are
infinitely many worlds in~$\mModel{M}$ that behave the same in
relevant respects, then we may be able to collect \emph{all} worlds in
finitely many (possibly infinite) ``classes'' of such worlds. In other
words, we should partition the set of worlds in the right way, i.e.,
in such a way that each partition contains infinitely many worlds, but
there are only finitely many partitions. Then we define a new
model~$\mModel{M^*}$ where the worlds are the partitions. Finitely
many partitions in the old model give us finitely many worlds in the
new model, i.e., a finite model. Let's call the partition a world~$w$
is in $[w]$. We'll want it to be the case that $\mSat{M}{!A}[w]$ iff
$\mSat{M^*}{!A}[{[w]}]$, since we want the new model to be a
countermodel to~$!A$ if the old one was. This requires that we define
the partition, as well as the accessibility relation of~$\mModel{M^*}$
in the right way.

To see how this would go, first imagine we have no accessibility
relation. $\mSat{M}{\Box !B}[w]$ iff for some $v \in W$,
$\mSat{M}{\Box !B}[v]$, and the same for $\mModel{M^*}$, except with
$[w]$ and $[v]$. As a first idea, let's say that two worlds $u$ and
$v$ are equivalent (belong to the same partition) if they agree on all
!!{propositional variable}s in~$\mModel{M}$, i.e., $\mSat{M}{p}[u]$
iff $\mSat{M}{p}[v]$. Let $V^*(p) = \Setabs{[w]}{\mSat{M}{p}[w]}$.
Our aim is to show that $\mSat{M}{!A}[w]$ iff
$\mSat{M^*}{!A}[{[w]}]$. Obviously, we'd prove this by induction: The
base case would be $!A \equiv p$. First suppose $\mSat{M}{p}[w]$. Then
$[w] \in V^*$ by definition, so $\mSat{M^*}{p}[{[w]}]$. Now suppose
that $\mSat{M^*}{p}[{[w]}]$. That means that $[w] \in V^*(p)$, i.e.,
for some $v$ equivalent to~$w$, $\mSat{M}{p}[v]$. But ``$w$ equivalent
to $v$'' means ``$w$ and $v$ make all the same !!{propositional variable}s
true,'' so $\mSat{M}{p}[w]$. Now for the inductive step, e.g., $!A
\ident \lnot !B$. Then $\mSat{M}{\lnot !B}[w]$ iff $\mSat/{M}{!B}[w]$
iff $\mSat/{M^*}{!B}[{[w]}]$ (by inductive hypothesis) iff
$\mSat{M^*}{\lnot !B}[{[w]}]$. Similarly for the other non-modal
operators. It also works for $\Box$: suppose $\mSat{M^*}{\Box
  !B}[{[w]}]$. That means that for every $[u]$,
$\mSat{M^*}{!B}[{[u]}]$. By inductive hypothesis, for every $u$,
$\mSat{M}{!B}[u]$. Consequently, $\mSat{M}{\Box !B}[w]$.

In the general case, where we have to also define the accessibility
relation for $\mModel{M^*}$, things are more complicated. We'll call a
model~$\mModel{M^*}$ a \emph{filtration} if its accessibility
relation~$R^*$ satisfies the conditions required to make the inductive
proof above go through. Then any filtration~$\mModel{M^*}$ will make
$!A$ true at $[w]$ iff $\mModel{M}$ makes $!A$ true at~$w$. However,
now we also have to show that there \emph{are} filtrations, i.e., we
can define~$R^*$ so that it satisfies the required conditions. In
order for this to work, however, we have to require that worlds $u$,
$v$ count as equivalent not just when they agree on all
!!{propositional variable}s, but on all sub-!!{formula}s
of~$!A$. Since $!A$ has only finitely many sub-!!{formula}s, this will
still guarantee that the filtration is finite. There is not just one
way to define a filtration, and in order to make sure that the
accessibility relation of the filtration satisfies the required
properties (e.g., reflexive, transitive, etc.) we have to be inventive
with the definition of~$R^*$.





% ==== OLP file: content/normal-modal-logic/filtrations/more-filtrations.tex ====
% Part: normal-modal-logic
% Chapter: filtrations
% Section: more-filtrations



\olfileid{nml}{fil}{acc}

\olsection{Filtrations and Properties of Accessibility}

As noted, filtrations of universal, serial, and reflexive models are
always also universal, serial, or reflexive. But not every filtration
of a symmetric or transitive model is symmetric or transitive,
respectively. In some cases, however, it is possible to define
filtrations so that this does hold. In order to do so, we proceed as
in the definition of the coarsest filtration, but add additional
conditions to the definition of~$R^*$.  Let $\Gamma$ be closed under
sub-!!{formula}s. Consider the relations~$C_i(u,v)$ in
\olref{tab:Cn-filtrations} between worlds $u$, $v$ in a model
$\mModel{M} =\tuple{W, R, V}$. We can define $R^*[u][v]$ on the basis
of combinations of these conditions. For instance, if we stipulate
that $R^*[u][v]$ iff the condition $C_1(u,v)$ holds, we get exactly
the coarsest filtration. If we stipulate $R^*[u][v]$ iff both
$C_1(u,v)$ and $C_2(u, v)$ hold, we get a different filtration. It is
``finer'' than the coarsest since fewer pairs of worlds satisfy
$C_1(u,v)$ and $C_2(u,v)$ than $C_1(u,v)$ alone.

\begin{table}[ht]
  \centering
  \begin{tabular}{|ll|}
    \hline
    \multirow{2}{*}{$C_1(u,v)$:}
    \iftag{prvBox}{&
      if $\Box!A \in \Gamma$ and $\mSat{M}{\Box!A}[u]$
      then $\mSat{M}{!A}[v]$;\iftag{prvDiamond}{ and}{}\\}{}
    \iftag{prvDiamond}{&
      if $\Diamond!A \in \Gamma$  and $\mSat{M}{!A}[v]$
      then $\mSat{M}{\Diamond!A}[u]$; \\}{}
    \hline
    \multirow{2}{*}{$C_2(u,v)$:} 
    \iftag{prvBox}{&
      if $\Box!A \in \Gamma$ and $\mSat{M}{\Box!A}[v]$
      then $\mSat{M}{!A}[u]$;\iftag{prvDiamond}{ and}{}\\}{}
    \iftag{prvDiamond}{&
      if $\Diamond!A \in \Gamma$ and $\mSat{M}{!A}[u]$
      then $\mSat{M}{\Diamond!A}[v]$; \\}{}
    \hline
    \multirow{2}{*}{$C_3(u,v)$:}
    \iftag{prvBox}{&
      if $\Box!A \in \Gamma$ and $\mSat{M}{\Box!A}[u]$
      then $\mSat{M}{\Box!A}[v]$;\iftag{prvDiamond}{ and}{}\\}{}
    \iftag{prvDiamond}{&
      if $\Diamond!A \in \Gamma$ and $\mSat{M}{\Diamond!A}[v]$
      then $\mSat{M}{\Diamond!A}[u]$; \\}{}
    \hline
    \multirow{2}{*}{$C_4(u,v)$:}
    \iftag{prvBox}{&
      if $\Box!A \in \Gamma$ and $\mSat{M}{\Box!A}[v]$
      then $\mSat{M}{\Box!A}[u]$;\iftag{prvDiamond}{ and}{}\\}{}
    \iftag{prvDiamond}{&
      if $\Diamond!A \in \Gamma$ and $\mSat{M}{\Diamond!A}[u]$
      then $\mSat{M}{\Diamond!A}[v]$; \\}{}
    \hline
    \end{tabular}
  \caption{Conditions on possible worlds for defining
    filtrations.}
  \ollabel{tab:Cn-filtrations}
\end{table}

\begin{thm}\ollabel{thm:more-filtrations}
  Let $\mModel{M} =\tuple{W,R,V}$ be a model, $\Gamma$ closed under
  sub-!!{formula}s. Let $W^*$ and $V^*$ be defined as in
  \olref[fil]{defn:filtration}. Then:
  \begin{enumerate}
  \item Suppose $R^*[u][v]$ if and only if $C_1(u, v) \land
    C_2(u,v)$. Then $R^*$ is symmetric, and $\mModel{M^*} =
    \tuple{W^*,R^*,V^*}$ is a filtration if $\mModel{M}$ is symmetric.
  \item Suppose $R^*[u][v]$ if and only if $C_1(u, v)
    \land C_3(u,v)$. Then $R^*$ is transitive, and
    $\mModel{M^*}=\tuple{W^*,R^*,V^*}$ is a filtration if $\mModel{M}$
    is transitive.
  \item Suppose $R^*[u][v]$ if and only if $C_1(u, v) \land C_2(u,v)
    \land C_3(u,v) \land C_4(u,v)$. Then $R^*$ is symmetric and
    transitive, and $\mModel{M^*}=\tuple{W^*,R^*,V^*}$ is a filtration
    if $\mModel{M}$ is symmetric and transitive.
  \item Suppose $R^*$ is defined as $R^*[u][v]$ if and only if $C_1(u,
    v) \land C_3(u,v) \land C_4(u,v)$. Then $R^*$ is transitive and
    euclidean, and $\mModel{M^*}=\tuple{W^*,R^*,V^*}$ is a filtration
    if $\mModel{M}$ is transitive and euclidean.
  \end{enumerate}
\end{thm}

\begin{proof}
  \begin{enumerate}
    \item It's immediate that $R^*$ is symmetric, since $C_1(u,v)
      \Leftrightarrow C_2(v,u)$ and $C_2(u,v) \Leftrightarrow
      C_1(v,u)$. So it's left to show that if $\mModel{M}$ is
      symmetric then $\mModel{M^*}$ is a filtration
      through~$\Gamma$. Condition $C_1(u,v)$ guarantees that
      \iftag{prvBox}{%
        \iftag{prvDiamond}{%
          \olref[fil]{defn:filtration-R2} and
          \olref[fil]{defn:filtration-R3} of
          \olref[fil]{defn:filtration} are}{%
          \olref[fil]{defn:filtration}\olref[fil]{defn:filtration-R2} is}}{%
        \olref[fil]{defn:filtration}\olref[fil]{defn:filtration-R3} is}
      satisfied. So we just have to verify
      \olref[fil]{defn:filtration}\olref[fil]{defn:filtration-R1},
      i.e., that $Ruv$ implies $R^*[u][v]$.

      So suppose $Ruv$. To show $R^*[u][v]$ we need to establish that
      $C_1(u,v)$ and $C_2(u,v)$. For $C_1$: \iftag{prvBox}{if $\Box!A
        \in\Gamma$ and $\mSat{M}{\Box!A}[u]$ then also
        $\mSat{M}{!A}[v]$ (since $Ruv$).\iftag{prvDiamond}{ Similarly,
        }{}}{}\iftag{prvDiamond}{if $\Diamond!A \in \Gamma$ and
        $\mSat{M}{!A}[v]$ then $\mSat{M}{\Diamond!A}[u]$ since
        $Ruv$.}{} For $C_2$: \iftag{prvBox}{if $\Box!A \in \Gamma$ and
        $\mSat{M}{\Box!A}[v]$ then $Ruv$ implies $Rvu$ by symmetry, so
        that $\mSat{M}{!A}[u]$.\iftag{prvDiamond}{ Similarly,
        }{}}{}\iftag{prvDiamond}{if $\Diamond!A \in\Gamma$ and
        $\mSat{M}{!A}[u]$ then $\mSat{M}{\Diamond!A}[v]$ (since $Rvu$
        by symmetry).}{}
    \item Exercise.
    \item Exercise.
    \item Exercise.
  \end{enumerate}
\end{proof}

\begin{prob}
  Complete the proof of \olref[nml][fil][acc]{thm:more-filtrations}.
\end{prob}


% ==== OLP file: content/normal-modal-logic/filtrations/preliminaries.tex ====
% Part: normal-modal-logic
% Chapter: filtrations
% Section: preliminaries



\olfileid{nml}{fil}{pre}

\olsection{Preliminaries}

Filtrations allow us to establish the decidability of our systems of
modal logic by showing that they have the \emph{finite model
  property}, i.e., that any !!{formula} that is true (false) in a
model is also true (false) in a \emph{finite} model.  Filtrations are
defined relative to sets of !!{formula}s which are closed under
subformulas.

\begin{defn}\ollabel{defn:modallyclosed}
  A set $\Gamma$ of !!{formula}s is \emph{closed under subformulas} if it
  contains every subformula of !!a{formula} in $\Gamma$. Further,
  $\Gamma$ is \emph{modally closed} if it is closed under subformulas
  and moreover $!A \in \Gamma$ implies $\Box!A,
  \Diamond!A \in \Gamma$. 
\end{defn}

For instance, given !!a{formula}~$!A$, the set of all its
sub-!!{formula}s is closed under sub-!!{formula}s. When we're defining
a filtration of a model through the set of sub-!!{formula}s of~$!A$,
it will have the property we're after: it makes $!A$ true (false) iff
the original model does.

The set of worlds of a filtration of~$\mModel{M}$ through~$\Gamma$ is
defined as the set of all equivalence classes of the following
equivalence relation.

\begin{defn}
Let $\mModel{M} =\tuple{W, R, V}$ and suppose $\Gamma$ is closed under
sub-!!{formula}s. Define a relation $\equiv$ on $W$ to hold of any two
worlds that make the same !!{formula}s from $\Gamma$ true, i.e.:
\[
u \equiv v \quad \text{if and only if } \quad 
\forall !A \in \Gamma : \mSat{M}{!A}[u] \Leftrightarrow \mSat{M}{!A}[v].
\]
The equivalence class~$[w]_\equiv$ of a world~$w$, or $[w]$ for short,
is the set of all worlds $\equiv$-equivalent to~$w$:
\[
[w] = \Setabs{v}{v \equiv w}.
\]
\end{defn}

\begin{prop}
  Given $\mModel{M}$ and $\Gamma$, $\equiv$ as defined above is an
  equivalence relation, i.e., it is reflexive, symmetric, and
  transitive.
\end{prop}

\begin{proof}
  The relation $\equiv$ is reflexive, since $w$ makes exactly the same
  !!{formula}s from~$\Gamma$ true as itself. It is symmetric since if
  $u$ makes the same !!{formula}s from~$\Gamma$ true as $v$, the same
  holds for $v$ and~$u$. It is also transitive, since if $u$ makes the
  same !!{formula}s from~$\Gamma$ true as~$v$, and $v$ as $w$, then
  $u$ makes the same !!{formula}s from~$\Gamma$ true as~$w$.
\end{proof}

The relation $\equiv$, like any equivalence relation, divides~$W$ into
\emph{partitions}, i.e., subsets of~$W$ which are pairwise disjoint,
and together cover all of~$W$. Every $w \in W$ is !!a{element} of one
of the partitions, namely of~$[w]$, since $w \equiv w$. So the
partitions $[w]$ cover all of~$W$. They are pairwise disjoint, for if
$u \in [w]$ and $u \in [v]$, then $u \equiv w$ and $u \equiv v$, and
by symmetry and transitivity, $w \equiv v$, and so $[w] = [v]$.


